\section{Conditional burst means}
\label{sec:cond-means}

The mean burst size and its various conditionings hide a small trap. It is
worth springing it deliberately.

\subsection{The overall expected burst size}

Recall that $V(t)=\mean{\wt}$, where $\wt$ is the number of particles
released by time $t$. Whether the catastrophe has taken place or not can be
represented by the truth of $\wt>0$, since the minimum burst size is $1$;
rupture has occurred with probability $1-I(t)$. Decomposing,
$$\mean{\wt} = \mean{\wt\mid\wt > 0}(1-I(t)) + {\mean{\wt \mid \wt = 0}}I(t),$$
and since ${\mean{\wt \mid \wt = 0}} = 0$,
$$\mean{\wt} = \mean{\wt\mid\wt > 0}(1-I(t)),$$
that is,
\begin{equation}
\mean{\wt\mid\wt > 0} = \frac{V(t)}{1-I(t)}.
\label{eq:VoverI}
\end{equation}
This can be read as ``mean burst size, given that a burst has happened by
$t$''. The overall expected burst size --- the expected size of an
arbitrarily generated burst --- is the late-time value of this expression:
\begin{equation}
\mean{\Kb}=\lim_{t\to\infty}\frac{V(t)}{1-I(t)}
=\frac{V_\infty}{1-b}
=\frac{\lambda-\mu}{\delta}+\frac{1}{1-b},
\label{eq:meanK}
\end{equation}
where the last form follows from \eqref{eq:Vinf}, and $\mean{\Kb}=a/(a-1)$
as already noted there. The distinction matters:
$V_\infty$ itself averages in the realisations that cleared internally and
released nothing, so $V_\infty$ is \emph{not} the mean burst size unless
$\mu=0$.

\subsection{The curious circularity, resolved}
\label{sec:circularity}

What is the expected burst size of a cell that has \emph{not yet} burst?
Decomposing $\mean{\Kb}$ on the event $\{\tau>t\}$,
\begin{align}
\mean{\Kb} &= \mean{\Kb\mid \tau > t} \pr{\tau > t} + \mean{\Kb\mid \tau < t}\pr{\tau<t}\nonumber\\
&= \mean{\Kb\mid\tau > t}\, I(t) + \frac{V(t)}{1 - I(t)} \lrb{1 - I(t)},
\end{align}
so that
\begin{equation}
\mean{\Kb\mid \tau > t}  = \frac{\mean{\Kb} - V(t)}{I(t)}
=\frac{V_\infty/(1-b) - V(t)}{I(t)}.
\label{eq:Kcond}
\end{equation}
So far, so good. The trap is what happens when one takes $t\to\infty$ in
\eqref{eq:Kcond}. A first attempt gives
$$
\lim_{t\to\infty}\mean{\Kb\mid \tau > t}
=\frac{V_\infty/(1-b)-V_\infty}{b},
$$
which, after a little algebra, appears to return $\mean{\Kb}$ itself --- a
curious circularity: the burst size of a cell that has waited an infinite
time to burst seems to equal the burst size of a typical cell, which makes
no sense, since conditioning on $\tau>t$ as $t\to\infty$ should select the
realisations that never burst at all.

The resolution is that \eqref{eq:Kcond} mixes two conditionings. Its
numerator $\mean{\Kb}-V(t)$ is the contribution to the burst size from
realisations that burst \emph{after} $t$; its denominator $I(t)$ is the
probability of not having burst by $t$, which includes the mass $b$ of
realisations that will never burst and whose burst size is zero. Written
consistently,
\begin{equation}
\mean{\Kb\,\mathbf 1\{\text{burst}\}\mid\tau>t}
=\frac{V_\infty-V(t)}{I(t)}
\xrightarrow[t\to\infty]{}\frac{V_\infty-V_\infty}{b}=0,
\end{equation}
which is correct and unmysterious: conditioning on ``has not burst yet''
converges to ``never bursts'', on which event the burst size is zero.

The genuinely interesting limit conditions on eventual burst as well:
\begin{equation}
\mean{\Kb\mid \tau>t,\ \text{burst eventually}}
=\frac{V_\infty-V(t)}{I(t)-b}.
\end{equation}
This is of the form $0/0$ as $t\to\infty$; by l'H\^opital's rule, using
$V'=\delta K$ and $I'=-\delta J$,
\begin{equation}
\lim_{t\to\infty}\mean{\Kb\mid \tau>t,\ \text{burst eventually}}
=\lim_{t\to\infty}\frac{K(t)}{J(t)}
=\frac{\langle X^2\rangle_{\rm QS}}{\langle X\rangle_{\rm QS}}
=1+\frac{2\lambda}{\delta}(1-b),
\label{eq:latemean}
\end{equation}
in agreement with \eqref{eq:sizebias}: the eventual burst size of a very
patient cell is the size-biased quasi-stationary mean. The circularity was
never in the mathematics; it was in the conditioning.
